
\medskip
\phantomsection\label{cor:sga3-VIB-6.2.6}%
\noindent\textbf{Corollary 6.2.6.---}%
\footnote{\label{note:sga3-VIB-59}Editor's note: we have inserted
this corollary here (cf.\ \SGARefLink{sga3:xii:6.1}{Expos\'e~XII, 6.1}), as it will be useful
later. We return in \SGARefLink{rem:sga3-VIB-6.3}{6.3} to the original text of
Expos\'e~VI\(_B\).}
\emph{Let \(k\) be a field and let \(G\) be a connected
algebraic \(k\)-group. Then
\(\underline{\operatorname{Centr}}_G(G)\) is representable by a
closed group subscheme \(Z\) of \(G\), and \(G/Z\) is an affine
algebraic \(k\)-group.}

\emph{Proof.}
The first assertion is, of course, contained in
\hyperref[cor:sga3-VIB-6.2.5]{6.2.5\textup{(i)}}, but we shall
see that it also follows from the proof of the second assertion.
Indeed, \(G\) acts by the adjoint representation on the
finite-dimensional \(k\)-vector spaces
\[
  V_n=\mathcal O_{G,e}/\mathfrak m_e^{\,n+1},
\]
where \(\mathfrak m_e\) is the maximal ideal of
\(\mathcal O_{G,e}\). Let \(K_n\) be the kernel of the morphism
\[
  \rho_n:G\longrightarrow\operatorname{GL}(V_n).
\]
By \hyperref[prop:sga3-VIA-5.4.1]{VI\(_A\), 5.4.1},
\(\rho_n\) induces a closed immersion
\[
  G/K_n\lhook\joinrel\longrightarrow\operatorname{GL}(V_n),
\]
so every \(G/K_n\) is affine. Since \(G\) is Noetherian, the
intersection \(K\) of the \(K_n\) is equal to one of the \(K_n\);
hence \(G/K\) is affine.

On the other hand, if \(Z\) denotes the center of \(G\), it is
clear that
\[
  Z\subset K.
\]
We shall prove the reverse inclusion. Let
\(\widehat{\mathcal O}_{G,e}\) be the completion of
\(\mathcal O_{G,e}\) for the \(\mathfrak m_e\)-adic topology,
and put
\[
  \widehat S=\operatorname{Spec}\widehat{\mathcal O}_{G,e},
  \qquad
  S=\operatorname{Spec}\mathcal O_{G,e}.
\]
Since \(\widehat S\to S\) is faithfully flat, and since the two
morphisms
\[
  K\times_k S\longrightarrow S,
  \qquad
  (g,x)\longmapsto gxg^{-1}
  \quad\text{and}\quad
  (g,x)\longmapsto x
\]
coincide after the base change \(\widehat S\to S\), they
coincide. Thus \(K\) acts trivially on \(\mathcal O_{G,e}\).

By \hyperref[ex:sga3-VIB-6.2.4-e]{6.2.4}\textup{(e)}, the subobject \(G^K\) of points of \(G\)
fixed by \(K\), which is precisely
\(\underline{\operatorname{Centr}}_G(K)\), is a closed
subscheme of \(G\). It is therefore defined by a quasi-coherent
ideal \(\mathcal I\) of \(\mathcal O_G\). Since \(G^K\) contains
\[
  S=\operatorname{Spec}\mathcal O_{G,e}
\]
and \(\mathcal I\) is of finite type (because \(G\) is
Noetherian), there is an open neighborhood \(U\) of \(e\) such
that
\[
  \mathcal I|_U=0.
\]
The subgroup
\[
  G^K=\underline{\operatorname{Centr}}_G(K)
\]
therefore contains \(U\), and hence also \(U\cdot U\), which is
equal to \(G\) because \(G\) is irreducible
(\hyperref[sga3:VIA:0.5]{VI\(_A\), 0.5}). Consequently,
\[
  \underline{\operatorname{Centr}}_G(K)=G.
\]
Thus \(K\subset Z\), and therefore \(Z=K\), as required.
